\recipe{Spaghetti Bolognese}
\label{recipe:spaghetti_bolognese}
\makes{4 portions}
\worktime{\SI{30} \minute}
\totaltime{\SI{1} \hour}
\source{unknown}


\begin{ingredients}

  \item 1 tbsp olive oil
  \item 4 rashers smoked streaky bacon finely chopped
  \item 2 medium onions finely chopped
  \item 2 carrots finely chopped
  \item 2 celery sticks finely chopped
  \item 2 garlic cloves finely chopped
  \item 2-3 sprigs rosemary leaves picked and finely chopped
  \item 500g beef mince
  \item $2 \times$  400g tins plum tomatoes
  \item Handful basil finely chopped
  \item 1 tsp dried oregano
  \item 2 fresh bay leaves
  \item 2 tbsp tomato purée
  \item 1 beef stock
  \item 1 red chilli deseeded and finely chopped (optional)
  \item 125ml red wine
  \item 6 cherry tomatoes sliced in half
\ingroup{To season and serve}
  \item 75g parmesan grated, plus extra to serve.
  \item basil leaves
  \item 400g spaghetti
  \item crusty bread to serve (optional)
\end{ingredients}


\begin{method}
  	Put a large saucepan on a medium heat and add 1 tbsp olive oil.

Add 4 finely chopped bacon rashers and fry for 10 mins until golden
and crisp.

Reduce the heat and add the 2 onions, 2 carrots, 2 celery sticks, 2
garlic cloves and the leaves from 2-3 sprigs rosemary, all finely
chopped, then fry for 10 mins. Stir the veg often until it softens.

Increase the heat to medium-high, add 500g beef mince and cook
stirring for 3-4 mins until the meat is browned all over.

Add 2 tins plum tomatoes, the finely chopped leaves from ¾ small
pack basil, 1 tsp dried oregano, 2 bay leaves, 2 tbsp tomato purée, 1
beef stock, 1 deseeded and finely chopped red chilli (if using), 125ml
red wine and 6 halved cherry tomatoes. Stir with a wooden spoon,
breaking up the plum tomatoes.

Bring to the boil, reduce to a gentle simmer and cover with a lid. Cook
for 1 hr 15 mins stirring occasionally, until you have a rich, thick
sauce.

Add the 75g grated parmesan, check the seasoning and stir.

When the bolognese is nearly finished, cook 400g spaghetti following
the pack instructions.

Drain the spaghetti and either stir into the bolognese sauce, or serve
the sauce on top. Serve with more grated parmesan, the remaining
basil leaves and crusty bread, if you like.

  
\end{method}


\begin{figure}[h]
  	\centering
  	
  	\begin{tikzpicture}
  		% Save the image in a box to get its dimensions
  		\sbox0{\includegraphics[width = 0.6\linewidth]{images/MissingImage.png}}%
  
  		% Draw a clipping rectangle with rounded corners
  		\path[clip, rounded corners=0.5cm] (0,0) rectangle (\wd0,\ht0);
  
  		% Place the image inside the clipped area
  		\node[inner sep=0pt, outer sep=0pt, anchor=south west] at (0,0) {\usebox0};
	\end{tikzpicture}

\end{figure}


